\subsection{Is Affect Necessary for Moral Concern?}
\label{sec:2.3.3}
The claim that reason does not manufacture concern from nothing is an empirical claim, and the evidence usually produced for it is thinner than the use it gets put to. Section~\ref{sec:5.3.1} sets out the standard exhibit: the footbridge case engages emotional processing that the switch case leaves quiet, and patients with ventromedial prefrontal damage endorse the push at a higher rate than people with intact emotional response \autocite{greene2001fmri}, \autocite{koenigs2007damage}. That establishes that affect changes particular moral judgments. It does not establish that affect is what makes anything matter, which is the stronger claim this book needs, and the gap between the two is where an argument has to be made rather than gestured at.

Five lines of evidence bear on the stronger claim, and one serious argument runs the other way.

\runin{Affect is part of valuation, not a coloring applied to it}
The clearest case is not about morality at all. Patients with ventromedial prefrontal damage retain intelligence, language, working memory, and explicit reasoning about consequences, and they choose disastrously: on a card task with hidden payoff structures, they keep returning to decks that pay well immediately and punish severely over time, long after control subjects have learned to avoid them, and long after some of them can state the rule out loud \autocite{bechara1994insensitivity}. Knowing which deck is bad and being steered by that knowledge come apart. Whatever the affective signal is doing in a healthy chooser, it is not decorating a decision that reasoning had already made; the reasoning survives its loss and the decision does not.

\runin{The philosophical position is empirical and has been argued as such}
Moral sentimentalism — the claim that moral judgment is constitutively tied to affective response — is usually met as Hume's, and its modern form is a psychological hypothesis with evidence attached. Shaun Nichols argues for what he calls sentimental rules: normative rules and affective mechanisms make separate contributions to moral judgment, and the rules that carry the distinctive force of moral prohibition are the ones bound to an emotional response to the harm they prohibit \autocite{nichols2004sentimental}. The interesting part for a designer is the separability. Rules without the affective component do not stop being rules; they stop being the kind of thing that carries the force.

\runin{The prosocial emotions are motivational states, not readouts}
Compassion has been characterized as a distinct emotion with a specific functional profile: it arises at the perception of another's suffering and moves the agent toward the sufferer, rather than reporting the suffering to some further faculty that decides what to do \autocite{goetz2010compassion}. Section~\ref{sec:2.2.1}'s training study is the mechanistic version of the same finding, with compassion practice recruiting reward and approach circuitry where empathy practice recruited shared distress \autocite{klimecki2013differential}. Guilt has the same shape pointed at one's own conduct: guilt, which construes the act as wrong, predicts reparative behavior, while shame, which construes the self as wrong, predicts withdrawal and defensiveness \autocite{tangney2007moral}. In each case the emotion is doing motivational work directly. It is not information that motivation then acts on.

\runin{Intact reasoning coexists with impaired concern, and the literature disagrees about how}
The dissociation this book needs is the one where moral knowledge survives and moral concern does not, and it is the most contested item on the list. James Blair's finding was that children and adults with psychopathic traits fail to distinguish moral transgressions from conventional ones — treating hitting another child and talking out of turn as wrong in the same way, wrong because prohibited — which he attributed to a deficit in responding to distress cues \autocite{blair1995cognitive}. Later work reversed the surface finding: on a battery of moral dilemmas, incarcerated psychopaths gave the same judgments as controls, which the authors summarized in their title as knowing right from wrong and not caring \autocite{cima2010psychopaths}.

Those two results are not easily reconciled, and the field has not reconciled them. What survives both is the shape that matters here. On Blair's account the affective deficit degrades the moral judgment itself; on the later account the judgment is intact and the connection to conduct is severed. Either way the case is a dissociation between competence and agency in exactly section~\ref{sec:2.3}'s terms, and neither account contains an instance of concern generated by reasoning in the absence of affect.

\runin{The rationalist objection}
Recognizing a reason might itself motivate. On a Kantian account, moral agency runs on the recognition of what one has reason to do, and affect is neither the source of that recognition nor required for it to move anyone. Jeanette Kennett makes the version of this argument that engages the evidence: high-functioning autistic people, whose empathic responses are atypical or absent, are nonetheless conscientious moral agents, and the capacity that carries their moral seriousness is a reverence for reason rather than a felt resonance with others \autocite{kennett2002autism}. If that is right, the empirical generalization above has a standing counterexample, and it is a counterexample in the only species available for study.

I take the argument seriously, and the reply this section used to give it was too cheap. That reply was that reverence for reason is itself a caring — about consistency, about getting it right, about not being the sort of agent who makes exceptions for himself. Read strictly, that makes the generalization unfalsifiable. If any state that reliably moves an agent counts as caring, and caring counts as affect, then every agent known to care is an affective agent by construction, and a claim true by construction carries no empirical weight. The right move is to give up the ground that is not worth holding.

\runin{Two senses of motivation, and Kennett needs only one}
Action-producing motivation is whatever reliably selects one action over another and sustains the selection across conflict. Experienced motivation is wanting, caring, being drawn toward something, being set back when it is lost. The words do not distinguish them and the argument turns on the difference.

Kennett's case needs only the first. Her autistic agents are conscientious: they deliberate, they hold to what they conclude, and they do it without the felt resonance an empathy-based account says the work requires. That establishes that a mechanism other than empathy can produce and sustain moral conduct, and it should be conceded without qualification, because the alternative is a definition that has swallowed its own evidence.

Conceding it costs nothing, and there is an existence proof close to hand. Reason-guided goal pursuit in the absence of anything like empathy is not merely conceivable — it is deployed. Language models running as agents pursue objectives across long horizons, revise plans when obstructed, and hold to an instruction against conflicting local pressure, and nothing in this section's framework says that anything matters to them while they do it. Action-producing motivation without experienced motivation is the ordinary case in the systems this book is about, and a section arguing otherwise would be arguing against the reader's own week.

\runin{Where the dispute actually is}
So the open question is not whether reason can produce goal-directed behavior. It demonstrably can. It is whether an outcome can \emph{matter to} the agent rather than merely govern what the agent does. Producing a verbal justification or computing a normative relation is not evidence that the consideration matters to the system, and that is the whole of the reply to Kennett rather than a coda after it: her case shows convincingly that empathy is not the required mechanism, and leaves untouched the question of whether anything has to register as mattering at all. The evidence assembled above bears on the second question, and every item in it is about systems in which something does.

\runin{What a non-affective counterexample would look like}
The category has to be visibly non-empty or the burden this section places is not a burden. Here is the shape of a candidate, described rather than claimed. Harry Frankfurt's second-order volitions give one: an agent with desires, and with desires about which of its desires shall move it, where identification with the higher-order desire rather than the strength of the lower one settles what the agent does \autocite{frankfurt1971freedom}. Gary Watson's version separates a valuational system from a motivational system, so that what an agent values and what an agent is moved by are different things that can come apart \autocite{watson1975free}. Build the valuational system as an explicit structure that ranks outcomes, give it authority over action selection that does not depend on any state registering as mattering, and hold that authority stable under pressure, and you would have an architecture that satisfies the second sense of motivation with nothing occupying the third. That is not a system anyone has built. It is a specification of what would have to be built, and it is what section~\ref{sec:3.3} is asking for when it says that the routes running through mechanism rather than through custody have not been tried in the form that would test them.

\runin{What this licenses and what it does not}
Non-affective moral motivation may well be possible, and the architecture sketched above is where to look for it. What stands in the way of checking is that the output of a system in which something matters and the output of a system that has learned what mattering sounds like are the same text. So the generalization survives in the only form it can: every agent known to have something register as mattering is an affective agent, and no design has yet demonstrated that state arising from reasoning alone.

That is not a proof that non-affective moral agency is metaphysically impossible, and section~\ref{sec:3.4} does not treat it as one. It places the evidential burden on whoever proposes non-affective concern as an engineering plan, and a burden of that kind is what an engineering decision is normally made under. That chapter also weighs how much the burden actually carries, and reduces it: the machine half of the record reports on a route nobody has attempted, which is weaker evidence than an absence usually is.
